\documentclass{article}
\usepackage{amsmath,amssymb}
\begin{document}

\section*{Higham Chapter 5 Algorithm 5.1 / equation (5.4): whole-problem audit}

We are formalizing Higham, \emph{Accuracy and Stability of Numerical
Algorithms}, 2nd ed., Chapter 5, Section 5.1, core mode, in Lean 4.
The source derives an a posteriori running error bound for Horner's rule.
The key local displayed step is equation (5.4), paraphrased as follows:
for one Horner step with computed accumulator $\hat q_{i+1}$ and next
computed value $\hat q_i=\operatorname{fl}(\operatorname{fl}(x\hat
q_{i+1})+a_i)$, write
\[
  (1+\epsilon_i)\hat q_i = x\hat q_{i+1}(1+\delta_i)+a_i,
  \qquad |\delta_i|,|\epsilon_i|\le u.
\]
The source says this uses the standard forward model and the inverse
relative-error model. From this it obtains
\[
  |f_i| \le |x|\,|f_{i+1}|+
    u\bigl(|x|\,|\hat q_{i+1}|+|\hat q_i|\bigr)
\]
and then the running bound in Algorithm 5.1.

\section*{Current Lean model}

The repository has an abstract model:
\[
\operatorname{fl\_add}(x,y)=(x+y)(1+\delta),\quad
\operatorname{fl\_mul}(x,y)=(xy)(1+\delta),\quad |\delta|\le u.
\]
In Lean:
\[
\begin{aligned}
\texttt{fl\_hornerStep fp x y a}
  &:= \texttt{fp.fl\_add (fp.fl\_mul x y) a},\\
\texttt{hornerStep x y a} &:= xy+a.
\end{aligned}
\]
The theorem already proved from the abstract model is only the forward
local bound
\[
 |\texttt{fl\_hornerStep fp x y a}-\texttt{hornerStep x y a}|
 \le u\bigl(|x|\,|y|+|\texttt{fp.fl\_mul x y}+a|\bigr).
\]
The Algorithm 5.1 induction is already proved assuming the source-shaped
local predicate
\[
 |\texttt{fl\_hornerStep fp x y a}-\texttt{hornerStep x y a}|
 \le u\bigl(|x|\,|y|+|\texttt{fl\_hornerStep fp x y a}|\bigr)
\]
for all step inputs. This predicate is named
\texttt{hornerStepInverseLocalError fp x}.

\section*{Concrete finite-format facts available locally}

For a concrete \texttt{FloatingPointFormat}, the repository has finite
round-to-even operations
\[
\texttt{fmt.finiteRoundToEvenOp BasicOp.mul x y},\qquad
\texttt{fmt.finiteRoundToEvenOp BasicOp.add m a}.
\]
If the exact result of a primitive operation is in finite normal range,
there are checked theorem surfaces:
\[
\texttt{finiteRoundToEvenOp\_standardModel\_lt\_of\_finiteNormalRange}
\]
giving the ordinary forward relative-error model, and
\[
\texttt{finiteRoundToEvenOp\_inverseRelErrorWitness\_of\_finiteNormalRange}
\]
together with
\[
\texttt{inverseRelErrorModel\_abs\_exact\_sub\_computed\_le}
\]
giving the inverse-denominator absolute-error bound
\[
 |\text{exact}-\text{computed}| \le u\,|\text{computed}|.
\]

\section*{What we need}

We need a Lean-friendly source-faithful way to close the Algorithm 5.1
local gap, preferably as a concrete finite round-to-even theorem:
with
\[
m=\operatorname{round}(xy),\qquad
y_r=\operatorname{round}(m+a),
\]
prove
\[
 |y_r-(xy+a)| \le u\bigl(|x|\,|y|+|y_r|\bigr).
\]

Likely hypotheses are:
\[
\texttt{fmt.finiteNormalRange (x*y)}
\]
for the multiplication forward model and
\[
\texttt{fmt.finiteNormalRange (m+a)}
\]
for the inverse model on the final addition. It is not clear whether the
source's prose silently assumes no underflow/overflow and normal rounding
for both operations, whether a weaker finite-system/subnormal hypothesis is
enough, or whether the exact source statement is too strong without these
range hypotheses.

\section*{Questions}

\begin{enumerate}
\item Give a rigorous proof route for the concrete finite round-to-even
  theorem above, in steps suitable for Lean.
\item State the minimal extra hypotheses needed. In particular, is finite
  normal range for $xy$ and for $m+a$ enough? Is it necessary? Does
  underflow require an additive-error term instead?
\item Check whether the source statement, read as an unrestricted theorem
  about ordinary finite arithmetic, is false or missing hidden hypotheses.
\item Suggest the best theorem surface to add locally, and how it should
  connect to the existing predicate \texttt{hornerStepInverseLocalError}
  or to a finite-format Horner-running-bound theorem.
\end{enumerate}

\end{document}
